\section{Agentes en la era de las API: de función universal a sistema ejecutable}

\subsection{La perspectiva de la función universal}
En condiciones de solo-API, el LLM puede considerarse como una «función universal condicionada por lenguaje natural». Tiene una gran generalidad a nivel textual, por lo que el peso de la ingeniería se traslada a «cómo incrustar esta función en un flujo ejecutable».

\begin{importantbox}{La API no es un nodo hoja, sino un controlador}
El docente enfatiza: en los sistemas de agentes modernos, la API del LLM no es solo una subllamada dentro de un programa fijo; con frecuencia determina directamente el flujo de ejecución, la selección de herramientas y la actualización de estados intermedios, desempeñando en realidad el papel de controlador.
\end{importantbox}

\begin{figure}[H]
\centering
\includegraphics[width=0.93\textwidth]{frame-03-abstraction.jpg}
\caption{Agentic Abstraction: abstracción unificada de modelo, herramientas y humano}
\end{figure}
\noindent\footnotesize Evidencia visual: subtítulos aproximadamente en 00:04:18--00:05:30, el docente comenta «API controls execution flow» y «memory stream / retrieval / action».\normalsize

\subsection{Flujo de memoria, recuperación y ciclo de reflexión}
La arquitectura básica de agente que presenta el docente puede descomponerse en cuatro componentes de un ciclo cerrado: observación (observation), flujo de memoria (memory stream), recuperación (retrieval) y ejecución (action).
Además, el sistema incorpora un mecanismo de «reflexión» (reflection) que hace que el modelo produzca resúmenes de alto nivel de las trayectorias pasadas, y luego reincorpora esos resúmenes a las decisiones posteriores.

\begin{knowledgebox}{Por qué la reflexión es tan importante en los agentes}
El fallo de un agente basado en llamadas a herramientas rara vez proviene de un único error de inferencia, sino de un sesgo acumulado a lo largo de una trayectoria de varios pasos.
El valor de la reflexión está en comprimir el «ruido de la trayectoria» en una señal de estrategia reutilizable, de modo que los pasos posteriores dejen de repetirse a ciegas.
\end{knowledgebox}

\begin{warningbox}{Las dos trampas más frecuentes de los sistemas de memoria}
\begin{itemize}
    \item Acumulación sin límite: el contexto se expande de manera continua y degrada la calidad de la recuperación;
    \item Fijación de errores: los resúmenes erróneos se citan de manera repetida y forman una «alucinación autorreforzada».
\end{itemize}
Por eso, la memoria no consiste en «almacenar cuanto más mejor», sino en que importa más la calidad de la compresión y la capacidad de verificación.
\end{warningbox}

\subsection{Resumen de la sección}
Lo esencial del agente en la era de las API no es «llamar a unas cuantas herramientas más», sino construir un ciclo de ejecución controlable: hacer del LLM un controlador y, al mismo tiempo, usar mecanismos de memoria y reflexión para reducir el error acumulado en cadenas largas.

